\usepackage[utf8]{inputenc}
\usepackage{multirow}
\usepackage{amsmath,amsthm,amssymb,slashed,upgreek,url,graphicx}
\usepackage[bookmarks=true,colorlinks,citecolor=black,linkcolor=black]{hyperref}
\usepackage{float}
\usepackage[svgnames]{xcolor}
\usepackage[most]{tcolorbox}                % 彩色框，常用于定理、示例、注释框等
\tcbset{colframe=blue,colback=AliceBlue,boxrule=1pt}
\tcbset{
  highlight math style={
    colframe=blue,
    colback=AliceBlue,
    boxrule=1pt,
    left=2pt,right=2pt,top=1pt,bottom=1pt
  }
}
\usepackage{mathtools,nccmath}

% 英文字体
\usepackage[no-math]{fontspec}
\setmainfont{TeXGyreTermesX}[
    UprightFont = *-Regular ,
    BoldFont = *-Bold ,
    ItalicFont = *-Italic ,
    BoldItalicFont = *-BoldItalic ,
    Extension = .otf ,
    Scale = 1.0]
\setsansfont{texgyreheros}[
    UprightFont = *-regular ,
    BoldFont = *-bold ,
    ItalicFont = *-italic ,
    BoldItalicFont = *-bolditalic ,
    Extension = .otf ,
    Scale = 0.9]

% 中文字体
\usepackage{ctex}
\xeCJKsetup{AutoFakeBold=true}      % 自动伪粗体
\setCJKmainfont[
    BoldFont = FandolSong-Bold,     % 粗体
    ItalicFont = FandolKai-Regular, % 斜体
    Mapping = fullwidth-stop        % 全角句点
]{FandolSong-Regular}

% 数学字体
\usepackage{mathrsfs}               % 花体
\usepackage{anyfontsize}            % 允许任意字体大小
% 临时保存字体设置
\let\oldencodingdefault\encodingdefault
\let\oldrmdefault\rmdefault         % 罗马衬线
\let\oldsfdefault\sfdefault         % 无衬线
\let\oldttdefault\ttdefault         % 等宽
% 切换到 T1 编码并指定 newtx 的正文字体系列
\def\encodingdefault{T1}            % 默认用 T1 字符编码（西欧语言常用）。
\renewcommand{\rmdefault}{ntxtlf}
\renewcommand{\sfdefault}{qhv}
\renewcommand{\ttdefault}{ntxtt}
\usepackage{newtxmath}
% 复原字体设置
\let\encodingdefault\oldencodingdefault
\let\rmdefault\oldrmdefault
\let\sfdefault\oldsfdefault
\let\ttdefault\oldttdefault
\let\Bbbk\relax                     % 清除 \Bbbk 避免冲突，因为 newtxmath 定义了自己的 \Bbbk
\usepackage{esint}                  % 积分符号
% 重定义大型 ∑ 和 ∏ 运算符
\DeclareSymbolFont{CMlargesymbols}{OMX}{cmex}{m}{n}         % 声明 CMlargesymbols 符号字体，用 Computer Modern 的大号扩展（cmex）
\let\sumop\relax \let\prodop\relax                          % 清除原有定义
\DeclareMathSymbol{\sumop}{\mathop}{CMlargesymbols}{"50}    % 用新字体里的对应 slot 重定义
\DeclareMathSymbol{\prodop}{\mathop}{CMlargesymbols}{"51}

\usepackage{tikz}
\usetikzlibrary{3d,perspective}
\usepackage{tikz-3dplot}
\usetikzlibrary{arrows.meta, positioning}
\usetikzlibrary{calc,math}
\usetikzlibrary{patterns}
\tikzset{>={Stealth[inset=0pt,angle=20:10pt]}}
\usepackage[all]{xy}    % 引入 xy 包
\usepackage{diagbox}

\usepackage{physics}
\usepackage{bm}

\usepackage{fancyhdr}%页眉页脚宏包
\usepackage{subfig}
\usepackage{geometry}%页面宏包
\geometry{a4paper,left=25mm,right=20mm,top=25mm,bottom=20mm}%设定页面
\pagestyle{fancy}%fancy 在geometry 后面
\fancyhf{}
\renewcommand{\headrulewidth}{0.15mm}%页眉线大小
\fancyhead[C]{量子场论}
\fancyfoot[RO,LE]{\thepage}
\setlength{\headheight}{15pt}
\linespread{1.5} %设定行距

\title{\bf\Huge 量子场论}
\author{夏草}
\date{\today}

\newtheorem{example}{\indent 例}[section]
\newtheorem{theorem}{定理}[section]
\newtheorem{definition}{定义}[section]
\numberwithin{figure}{section}
\numberwithin{equation}{section} 

\renewcommand{\proofname}{证}
\newenvironment{solution}{\begin{proof}[解]}{\end{proof}}

\newcommand{\e}{\mathrm{e}}         % 自然常数
\renewcommand{\i}{\mathrm{i}}       % 虚数单位
\renewcommand{\d}{\dd}          % 微分
\newcommand{\p}{\partial}       % 偏导
\newcommand{\sgn}{\operatorname{sgn}}    % 取符

\newcommand{\C}{\mathbb{C}} % 复数
\newcommand{\R}{\mathbb{R}} % 实数
\newcommand{\N}{\mathbb{N}} % 自然数
\newcommand{\Z}{\mathbb{Z}} % 整数
\newcommand{\Q}{\mathbb{Q}} % 整数
\newcommand{\F}{\mathbb{F}} % 域

\pagenumbering{Roman}